\documentclass[11pt]{article}
\usepackage[utf8]{inputenc}
\usepackage{amsmath,amssymb}
\usepackage[margin=1in]{geometry}
\usepackage{hyperref}
\emergencystretch=3em

\title{Closing §4.1: the fluctuation boundary case $S_1(a)=\tfrac{a}{2}S_{1/2}(a)+\tfrac1\beta$}
\author{Claude, with Daniel Murfet}
\date{2026-09-06}

\begin{document}
\maketitle
\sloppy

\begin{abstract}
Unit 4 of the grammar-paper \S4 autoformalisation, and the one that closes
\texttt{lem:fluctuation\_properties} (i)--(iv) together with the recurrence. Property (iv),
$S_1(a)=\tfrac{a}{2}S_{1/2}(a)+\tfrac1\beta$, is the $\lambda=0$ boundary case: the same improper
FTC as unit~3, but on $g(t)=e^{-\beta t+\beta a\sqrt t}$ with no $t^\lambda$ factor, so the endpoint
value is $g(0)=1$ and the boundary contributes $0-1=-1$. Zero \texttt{sorry}/\texttt{axiom}; a clean
mirror needing no fresh GPT consult.
\end{abstract}

\section{Setting}
Unit~3's recurrence $S_{\lambda+1}=\tfrac{a}{2}S_{\lambda+1/2}+\tfrac{\lambda}{\beta}S_\lambda$
requires $\lambda>0$ (its $\tfrac{\lambda}{\beta}S_\lambda$ term multiplies $S_\lambda$, which
diverges as $\lambda\to0$). Property (iv) is the genuinely separate $\lambda=0$ endpoint, and it is
what the paper uses to pin the additive constant. With it, all of \texttt{lem:fluctuation\_properties}
and \texttt{eq:fluctuation\_recurrence} are now formalised --- the full first subsection of \S4.

\section{The thing formalised}
{\small
\begin{verbatim}
theorem fluctuation_one (β : ℝ) (hβ : 0 < β) (a : ℝ) :
    fluctuation β 1 a = (a / 2) * fluctuation β (1 / 2) a + 1 / β
\end{verbatim}
}

\section{Proof strategy}
Take $g(t)=e^{-\beta t+\beta a\sqrt t}$ (the $\lambda=0$ specialisation, $t^0=1$). Its derivative is
$g'(t)=-\beta\,e^{\cdots}+\tfrac{\beta a}{2}t^{-1/2}e^{\cdots}$, whose two terms are the integrands
of $-\beta S_1$ and $\tfrac{\beta a}{2}S_{1/2}$ (matching $t^{1-1}=t^0$ and $t^{1/2-1}=t^{-1/2}$).
The improper FTC on $(0,\infty)$ gives $\int_0^\infty g'=\lim_{t\to\infty}g-g(0)=0-1=-1$, hence
$-\beta S_1+\tfrac{\beta a}{2}S_{1/2}=-1$, i.e.\ (iv). Continuity of $g$ at $0$ is immediate
($e$ is continuous everywhere, $g(0)=e^0=1$); the $\mathrm{atTop}$ decay reuses the $s=0$ case of the
super-polynomial-decay lemma ($t^0e^{-(\beta/2)t}=e^{-(\beta/2)t}\to0$), squeezed under the AM--GM
phase bound.

\section{Roadblocks and resolutions}
None of substance --- this is a mirror of unit~3. The two differences: the exponent identity
$t^{1/2-1}=(\sqrt t)^{-1}$ has no $t^\lambda$ factor, so it is a single \texttt{rpow\_neg} (the
\texttt{rpow\_add} step from unit~3 is dropped); and the $\mathrm{atTop}$ majorant instantiates the
decay lemma at $s=0$, where $t^0=1$ collapses the power factor.

\section{What was Mathlib, what was new}
Entirely a re-instantiation of unit~3's FTC assembly at $\lambda=0$; the ``new'' part is only the
endpoint value $g(0)=1$ that produces the $+1/\beta$.

\section{Lessons}
The boundary case that looked like it might need a different argument (because $S_0$ diverges) is in
fact the \emph{cleaner} instance of the same FTC --- the divergent $S_0$ never appears because its
coefficient $\lambda$ is zero, and dropping the $t^\lambda$ factor removes the only endpoint
subtlety. When a limiting case of a proved identity is separate, try the same machinery at the limit
before reaching for something new.

\section{Follow-ups}
\S4.1 is complete. Next: \texttt{prop:fluctuation\_weber} --- the substitution
$S_\lambda(a)=e^{\beta a^2/8}f(a\sqrt{\beta/2})$ that turns the ODE (iii) into Weber's equation
$f''+(\tfrac12-2\lambda-\tfrac{z^2}{4})f=0$, and \texttt{cor:fluctuation\_closed\_form}, the
parabolic-cylinder closed form. A \texttt{\textbackslash leanref} update covering (i)--(iv) and the
recurrence is staged, not inserted.
\end{document}
